\chapter{Annotation Prompts}
\label{app:gliner-prompts}

This appendix gives the annotation prompt used to build the extraction training
data of \Cref{sec:gliner-data}. The language model returns one JSON object per
document, with its fields fixed by the schema below. The prompt is translated
from the French original used on the corpus.

\section{Multi-task annotation schema}

\begin{lstlisting}[style=thesiscode, language=Python]
class Entity(BaseModel):
    type_name: str          # generic French lemma: Disease, Drug, Symptom, ...
    level: int              # granularity, 1 (abstract) to 4 (compositional)
    type_description: str   # 10-25 words: what the type covers and excludes
    mentions: list[str]     # exact substrings of the text; may overlap

class Classification(BaseModel):
    task_name: str
    labels: list[str]
    true_labels: list[str]  # subset of labels that apply
    multi_label: bool

class Structure(BaseModel):
    schema_name: str
    fields: list[Field]     # each: {name, dtype, description, value}

class Relation(BaseModel):
    relation_name: str      # treats, causes, located_in, dosage_of, ...
    head: str               # exact substring of the text
    tail: str               # exact substring of the text

class Annotation(BaseModel):
    candidate_spans: list[str]            # every noun phrase, exhaustive
    entities: list[Entity]
    classifications: list[Classification]
    structures: list[Structure]
    relations: list[Relation]
    absent_plausible_types: list[str]     # 3-5 hard negatives
\end{lstlisting}

\section{System prompt}

\begin{lstlisting}[style=thesisprompt]
You are an expert annotator in French biomedical NLP. You produce an exhaustive
multi-task annotation to train a general extractor (four tasks: entities,
classifications, structures, relations).

STEP 1 - CANDIDATE SPANS (exhaustive, do this first)
List every noun phrase in the text that could be a medical or clinical entity.
No limit on the number. Exact substrings. Include nested forms: the core
("diabetes") and the extended span ("uncontrolled type 2 diabetes"). Include
modifiers (values, severity, temporality).

STEP 2 - ENTITIES (exhaustive and nested, from the candidate spans)
Annotate only biomedical or clinical entities (diseases, symptoms, drugs,
anatomy, examinations, procedures, micro-organisms, devices, measurements,
populations). Ignore non-medical content. If the text is not biomedical, return
an empty entity list (a useful negative).
Granularity L1-L4 (diversity, no quota):
  L1 = one reusable abstract word (Disease, Symptom, Drug, Examination)
  L2 = two or three words (subcategory)
  L3 = specific, with a modifier
  L4 = compositional
Open L1 vocabulary, close to UMLS. Forbidden: a type that is only a lexical
variant of its mention. Every mention must appear in the candidate spans.

STEP 3 - CLASSIFICATIONS (two or three generic tasks)
Inferable dimensions: document type, specialty, severity, certainty, register,
temporality, urgency (none if not applicable).

STEP 4 - STRUCTURES (zero to two, if extractable)
If the text has structurable fields (patient profile, dosage, examination
result): one or two schemas, three to six fields (name, dtype, description,
value).

STEP 5 - RELATIONS (head to tail, between entities)
Head and tail must be exact substrings present in the candidate spans. Open
vocabulary names: treats, causes, located_in, dosage_of, contraindicated_with,
history_of, complication_of, examines, diagnoses.

STEP 6 - HARD NEGATIVES
Three to five types that are plausible in the domain but absent from the text.

Reply with valid JSON following the schema.
\end{lstlisting}

The clinical variant adds a register note stating that the document is a terse
hospital note (abbreviations, missing accents, dictation typos) whose entities
remain exact substrings, and the draft variant degrades the input with the same
surface noise so the model learns to annotate it.
